% !TEX TS-program = pdflatex
\documentclass[10pt,a4paper]{altacv}

% Ajustement des marges pour 1 seule colonne
\geometry{left=1.5cm,right=1.5cm,top=1.2cm,bottom=0.5cm,footskip=1\baselineskip}

\usepackage[T1]{fontenc}
\usepackage[utf8]{inputenc}
\usepackage[default]{lato}
\usepackage{microtype}
\DisableLigatures{encoding=*, family=*} % ATS-fix: Disables fi/fl/ffi ligatures in pdfLaTeX

\definecolor{Navy}{HTML}{0A2540}
\definecolor{SlateGrey}{HTML}{2E2E2E}
\definecolor{LightGrey}{HTML}{666666}
\colorlet{heading}{Navy}
\colorlet{accent}{Navy}
\colorlet{emphasis}{SlateGrey}
\colorlet{body}{LightGrey}

\renewcommand{\itemmarker}{{\small\textbullet}}
\renewcommand{\ratingmarker}{\faCircle}

\usepackage[hidelinks]{hyperref}
\hypersetup{
  pdftitle={Curriculum Vitae -- Achraf Izelmadene},
  pdfauthor={Achraf Izelmadene},
  pdfsubject={Ingénieur Génie Électrique et Systèmes Intelligents},
  pdfkeywords={génie électrique, IoT, automatisation, AutoML, robotique, PFE, stage}
}
\photo{2.2cm}{assets/imageee.png}

\renewcommand{\cvsection}[2][]{%
  \medskip%
  {\color{heading}\LARGE\bfseries #2 \hspace{0.2em}\leaders\vrule height 4pt depth -2pt\hfill\mbox{}}\par\smallskip
}
\renewcommand{\divider}{\textcolor{body!30}{\hdashrule{\linewidth}{0.6pt}{0.5ex}}\smallskip}

\renewcommand{\cvevent}[4]{%
  {\large\bfseries\color{accent}#1\par}
  \smallskip\normalsize
  \ifstrequal{#2}{}{}{%
    \textbf{\color{emphasis}#2\ifstrequal{#4}{}{}{ - #4}}%
  }%
  \hfill
  \ifstrequal{#3}{}{}{{\small#3}}%  % Removed calendar icon
  \par\smallskip
}

\makeatletter
\renewcommand{\makecvheader}{%
  \begingroup
    \ifdef{\@photodiameter}{%
    \begin{minipage}{\@photodiameter}
    \if@normalphoto
      \includegraphics[width=\linewidth]{\@photo}
    \else
      \begin{tikzpicture}
        \clip (0,0) circle (0.5\linewidth);
        \node at (0,-0.1\linewidth) {\includegraphics[width=1.05\linewidth]{\@photo}};
      \end{tikzpicture}
    \fi%
    \end{minipage}\hspace{1.5em}%
    \begin{minipage}{\dimexpr\linewidth-\@photodiameter-2em}}{}%
    \raggedright\color{emphasis}%
    {\Huge\bfseries\MakeUppercase{\@name}\par}
    \medskip
    {\large\bfseries\color{accent}\@tagline\par}
    \medskip
    {\footnotesize\bfseries\@personalinfo\par}
    \ifdef{\@photodiameter}{%
    \end{minipage}\par}{}%
  \endgroup\medskip
}
\makeatother

\begin{document}

\name{Izelmadene Achraf}
\tagline{Élève Ingénieur en Génie Électrique et Systèmes Intelligents}
\personalinfo{
  \email{\href{mailto:achraf.izelmadene@gmail.com}{achraf.izelmadene@gmail.com}}
  \phone{+212699098736}
  \linkedin{\href{https://www.linkedin.com/in/achraf-izelmadene/}{achraf-izelmadene}}
  \location{Marrakech, Maroc}
}

\makecvheader
\cvsection{PROFIL}
Élève ingénieur en 5\textsuperscript{ème} année de Génie Électrique et Systèmes Intelligents, je suis à la recherche d'un stage de fin d'études (PFE) à partir du 1\textsuperscript{er} février 2027. Animé par les défis technologiques du secteur ferroviaire, je souhaite mettre ma rigueur et mon expertise en ingénierie électrique au service de la conception et de la sécurisation de vos systèmes de signalisation.

\cvsection{EXPÉRIENCE}
\cvevent{Stagiaire Bureau d'Études}{Core Engineering Solutions}{juil. - août 2026}{Berrchid, Casablanca}
\begin{itemize}
\item Dimensionnement d'un poste HTA/BT et intégration d'une installation solaire PV en autoconsommation.
\item Réalisation de schémas unifilaires, bilans de puissance et notes de calcul (AutoCAD, Caneco BT, PVsyst).
\end{itemize}
\cvevent{Stagiaire en Unité de Production}{Ménara Préfa}{juil. 2024}{Marrakech}
\begin{itemize}
\item Migration et automatisation du processus de production vers un automate Siemens S7 sous TIA Portal.
\item Suivi des indicateurs d'équipements et observation active des opérations de maintenance préventive/corrective.
\end{itemize}

\cvsection{PROJETS}
{\large\bfseries\color{accent} Precision Agriculture IoT Ecosystem}\\[0.05cm]
\begin{itemize}
  \item Conception de nœuds Arduino Nano + LoRa, décision embarquée ESP32 (logique floue), et dashboard.
  \item Modèle d'IA TFLite Edge déployé sur smartphone pour la détection de maladies par vision embarquée.
\end{itemize}

{\large\bfseries\color{accent} Maintenance Prédictive 4.0 par AutoML}\\[0.05cm]
\begin{itemize}
  \item Comparaison d'outils AutoML TPOT vs. Auto-WEKA sur des données IoT (vibrations, température).
  \item Validation croisée temporelle, SMOTE, pipelines ML automatisés et étude CND par analyse d'huile.
\end{itemize}

{\large\bfseries\color{accent} Détection d'anomalies Smartgrid (Clustering)}\\[0.05cm]
\begin{itemize}
  \item Évaluation approfondie d'algorithmes (KMeans, DBSCAN) sur des données IoT de réseaux électriques.
  \item Développement d'un dashboard interactif dédié au monitoring continu et à la détection d'anomalies.
\end{itemize}

{\large\bfseries\color{accent} Commande d'un Gradateur par Arduino}\\[0.05cm]
\begin{itemize}
  \item Modélisation précise sous Simulink et interfaçage Hardware-in-the-Loop (HIL) avec une carte Arduino.
  \item Contrôle et régulation de la puissance électrique en temps réel via le pilotage de l'angle.
\end{itemize}


\cvsection{FORMATION}
\cvevent{Étudiant en cycle d'ingénieur en Génie Électrique et Systèmes Intelligents}{École Marocaine des Sciences de l'Ingénieur}{2025 - Présent}{Marrakech}
\small Réseaux et Sécurité électrique, Ingénierie Système, Automatisme Industriel, Électronique de puissance, Supervision \& SCADA...

\cvevent{Licence SMP, option Automatique}{Faculté des Sciences Semlalia}{2020 - 2024}{Marrakech}
\small Électrostatique \& Électromagnétisme, Électronique, Traitement du signal, Asservissement, Capteurs et Instrumentation...
%\cvevent{Baccalauréat Scientifique, option Sciences Physiques et Chimiques}{Lycée Maghreb Al Arabi}{2019 -- 2020}{Marrakech}

\cvsection{COMPÉTENCES}
\begin{tabular}{@{}ll@{}}
\textbf{Ingénierie \& CAO :} & Caneco BT, PVsyst, AutoCAD, SysML, MATLAB, Simulink \\[1pt]
\textbf{Programmation :} & Python, C/C++, JavaScript, HTML/CSS \\[1pt]
\textbf{Outils \& Plateformes :} & 3DEXPERIENCE, Arduino, Git, Docker, PyBullet \\[1pt]
\textbf{IA \& MLOps :} & Machine Learning, Scikit-learn, OpenCV, MediaPipe, AutoML, MLOps \\
\end{tabular}

\cvsection{CERTIFICATIONS}
\begin{itemize}[leftmargin=*, label=\textbullet, itemsep=2pt]
  \item \textbf{Foundations of IoT Systems and Industrial Automation} \hfill {\color{emphasis}LearnQuest (Coursera)}
  \item \textbf{Sensors and Sensor Circuit Design} \hfill {\color{emphasis}University of Colorado Boulder (Coursera)}
  \item \textbf{Certified 3DEXPERIENCE Collaborative Industry Innovator} \hfill {\color{emphasis}Dassault Systèmes}
  \item \textbf{MATLAB / Simulink Onramp} \hfill {\color{emphasis}MathWorks}
\end{itemize}

\cvsection{LANGUES \& DIVERS}
\textbf{Arabe :} Langue maternelle \hfill \textbf{Français :} DELF B2 \hfill \textbf{Anglais :} EFSET B2\\[2pt]
\textbf{Prix :} 1\textsuperscript{er} Hack The Intelligence - AISEC Summit 2026 \hfill 2\textsuperscript{ème} Vision AI Competition - EMSI 2026\\[2pt]
\textbf{Permis de conduire :} B \hfill \textbf{Loisirs :} Lecture, Échecs, Électronique DIY


\end{document}
